%======================================================================
%	
%	 __  __       _         ______ _ _      
%	|  \/  |     (_)       |  ____(_) |     
%	| \  / | __ _ _ _ __   | |__   _| | ___ 
%	| |\/| |/ _` | | '_ \  |  __| | | |/ _ \
%	| |  | | (_| | | | | | | |    | | |  __/
%	|_|  |_|\__,_|_|_| |_| |_|    |_|_|\___|
%
%----------------------------------------------------------------------
% Descripton : Main File to compose the resulting PDF
%======================================================================
%\pdfobjcompresslevel=0
\pdfminorversion=7
\pdfinclusioncopyfonts=1
\documentclass[
				ngerman             % neue deutsche Rechtschreibung
				,a4paper            % Papiergrösse
				,twoside            % Zweiseitiger Druck (rechts/links)
				% ,10pt             % Schriftgrösse
				,11pt
				% ,12pt
				,pdftex
				%  ,disable         % Todo-Markierungen auschalten
				,notitlepage		% Thanks abstract... preventing it from removing pageNumber
				]{report}

%————————————————————————————————————————————————————————————————————————————
%					Import aller Konfigurationen
%————————————————————————————————————————————————————————————————————————————
\usepackage{./config/DHBW/bericht}
\usepackage{./config/PaketeProjektarbeit} % <- spezifische Pakete für die Projektarbeit
\usepackage{./config/Pakete}
\usepackage{./config/Befehle}
\usepackage{./config/Layout}
\usepackage{./config/Styles}
%Wenn PDF/A-1B gewünscht .. Achtung am besten nochmal mit veraPDF testen!
\usepackage{./config/PDFStandard}

%————————————————————————————————————————————————————————————————————————————
%					Variablen deklarieren
%————————————————————————————————————————————————————————————————————————————
%======================================================================
%	
%	__      __        _       _     _            
%	\ \    / /       (_)     | |   | |           
%	 \ \  / /_ _ _ __ _  __ _| |__ | | ___ _ __  
%	  \ \/ / _` | '__| |/ _` | '_ \| |/ _ \ '_ \ 
%	   \  / (_| | |  | | (_| | |_) | |  __/ | | |
%	    \/ \__,_|_|  |_|\__,_|_.__/|_|\___|_| |_|
%
%----------------------------------------------------------------------
% Descripton : Variables
%======================================================================

\newboolean{DEBUG}
\newboolean{BERICHTE}
\newboolean{SPERRVERMERK}
\newboolean{CHANGELOG}
\newboolean{VORWORT}
\newboolean{ABSTRACT}
\newboolean{DANKSAGUNG}
\newboolean{WATERMARK}
\newboolean{INDEX}
\newboolean{TRENNSTRICHE}
\newboolean{ERKLAERUNG}
\newboolean{DOKUMENTENVERZEICHNIS}
\newboolean{DATEIVERZEICHNIS}
\newboolean{LITERATURVERWEIS}
\newboolean{DIGITALE-UNTERSCHRIFT}
\newboolean{DIGITAL-SIGN-AREA}

%————————————————————————————————————————————————————————————————————————————
%					Informationen ausfüllen
%————————————————————————————————————————————————————————————————————————————
%======================================================================
%	
%	_____        __                           _   _             
%  |_   _|      / _|                         | | (_)            
%	 | |  _ __ | |_ ___  _ __ _ __ ___   __ _| |_ _  ___  _ __  
%    | | | '_ \|  _/ _ \| '__| '_ ` _ \ / _` | __| |/ _ \| '_ \ 
%	_| |_| | | | || (_) | |  | | | | | | (_| | |_| | (_) | | | |
%  |_____|_| |_|_| \___/|_|  |_| |_| |_|\__,_|\__|_|\___/|_| |_|
%
%----------------------------------------------------------------------
% Descripton : File with all Information about the Document
%======================================================================

\newcommand{\Autor}{Max Mustermann}
\newcommand{\MatrikelNummer}{4711}
\newcommand{\Kursbezeichnung}{tinf17b3}

\newcommand{\FirmenLogoDeckblatt}{\includegraphics[width=2cm]{./config/layout/bonkpartners}}

% Falls es kein Firmenlogo gibt:
%  \newcommand{\FirmenLogoDeckblatt}{}

\newcommand{\BetreuerDHBW}{Mirko Dostmann}

%%%%%%%%%%%%%%%%%%%%%%%%%%%%%%%%%%%%%%%%%%%%%%%%%%%%%%%%%%%%%%%%%%%%%%%%%%%%%%%%%%%%%

% Wird auf dem Deckblatt und in der Erklärung benutzt:
%\newcommand{\Was}{Projekt-/Studien-/Bachelorarbeit}
%\newcommand{\Was}{Projektrarbeit}
%\newcommand{\Was}{Studienarbeit}
%\newcommand{\Was}{Bachleorarbeit}
\newcommand{\Was}{Programmentwurf}

%%%%%%%%%%%%%%%%%%%%%%%%%%%%%%%%%%%%%%%%%%%%%%%%%%%%%%%%%%%%%%%%%%%%%%%%%%%%%%%%%%%%%

\newcommand{\Titel}{\LaTeX-Vorlage für diverse Ausarbeitungen .. oder so ähnlich}
\newcommand{\AbgabeDatum}{1. April 2090}

\newcommand{\Dauer}{5. \& 6. Semester}

% \newcommand{\Abschluss}{Bachelor of Engineering}
\newcommand{\Abschluss}{Bachelor of Science}

\newcommand{\Studiengang}{Informatik / Informationstechnik}
% \newcommand{\Studiengang}{Informatik / Angewandte Informatik}
%======================================================================
%	 _    _                                 __ 
%	| |  | |                               / _|
%	| |__| |_   _ _ __   ___ _ __ _ __ ___| |_ 
%	|  __  | | | | '_ \ / _ \ '__| '__/ _ \  _|
%	| |  | | |_| | |_) |  __/ |  | | |  __/ |  
%	|_|  |_|\__, | .__/ \___|_|  |_|  \___|_|  
%		     __/ | |                           
%		    |___/|_|                           
%
%----------------------------------------------------------------------
% Descripton : Hyperref File for setting PDF Information
%======================================================================

\hypersetup{
	bookmarks=true,
	unicode=false,
	pdftoolbar=true,
	pdfmenubar=true,
	pdffitwindow=false,
	pdfstartview={FitH},
	pdftitle={\Titel},
	pdfauthor={\Autor},
	pdfsubject={\Was},
	pdfcreator={\Autor},
	pdfkeywords={DHBW},
	pdfborderstyle={/S/U/W 1},
	pdfnewwindow=true,
	citebordercolor=LimeGreen,
	filebordercolor=Magenta,
	linkbordercolor=Orange,
	menubordercolor=Orange,
	urlbordercolor=Fuchsia,
	runbordercolor=Orange
}


%————————————————————————————————————————————————————————————————————————————
%					Sonstige Einstellungen
%————————————————————————————————————————————————————————————————————————————
\bibliography{./directories/bibliography}
\loadglsentries{./directories/glossary}
\makeglossaries

%————————————————————————————————————————————————————————————————————————————
%					Einstellungen
%————————————————————————————————————————————————————————————————————————————
% Debug auf true setzen, hierbei wird der Changelog hinzugefügt, und keine Doppelseitenumbrüche gemacht (Zum drucken -> false)
% Bitte zum drucken oben in der Präambel auch auf Doppelseite schalten!
\setboolean{DEBUG}{true}
% Soll eine Erklärung eigefügt werden ?
\setboolean{ERKLAERUNG}{false}
%Soll das Abstract eingefügt werden ?
\setboolean{ABSTRACT}{false}
%Soll ne Trennlinie im Inhaltsverzeichnis eingefügt werden, vor Verzeichnissen, dem Inhalt und dem Anhang
\setboolean{TRENNSTRICHE}{true}
%Setzt ein Watermark
\setboolean{WATERMARK}{false}
%Dokumentenverzeichnis generieren ?

%————————————————————————————————————————————————————————————————————————————
%					Watermark
%————————————————————————————————————————————————————————————————————————————
\ifthenelse{\boolean{WATERMARK}}{\usepackage{draftwatermark}}{}
\ifthenelse{\boolean{WATERMARK}}{\SetWatermarkText{Entwurf}}{}

%————————————————————————————————————————————————————————————————————————————
%									Anfang Dokument
%————————————————————————————————————————————————————————————————————————————
\begin{document}
	% Wer keine Seitenzahlen braucht bis zum TOC einfach das Pagenumbering auf gobble und an der Stelle an der die Seitenzählung beginnen soll roman einfügen
	\pagenumbering{gobble}
	\pagestyle{plain}
	
	%—————————————————————————————————————————————————————————————————————————
	%								Deckseite
	%—————————————————————————————————————————————————————————————————————————
	\begin{titlepage}
		%%%%%%%%%%%%%%%%%%%%%%%%%%%%%%%%%%%%%%%%%%%%%%%%%%%%%%%%%%%%%%%%%%%%%%%%%%%%%%%
%% Descr:       Vorlage für Berichte der DHBW-Karlsruhe, Titlepage
%% Author:      Prof. Dr. Jürgen Vollmer, vollmer@dhbw-karlsruhe.de
%% $Id: erklaerung.tex,v 1.11 2020/03/13 14:24:42 vollmer Exp $
%% -*- coding: utf-8 -*-
%%%%%%%%%%%%%%%%%%%%%%%%%%%%%%%%%%%%%%%%%%%%%%%%%%%%%%%%%%%%%%%%%%%%%%%%%%%%%%%

\begin{center}
	\vspace*{-2cm}
	\FirmenLogoDeckblatt\hfill\includegraphics[width=4cm]{./config/DHBW/dhbw-logo.png}\\[2cm]
	{\Huge \Titel}\\[1cm]
	{\Huge\scshape \Was}\\[1cm]
	{\large der Vorlesung \hk{\textbf{A}dvanced \textbf{S}oftware \textbf{E}ngineering}}\\[0.5cm]
	{\large an der}\\[0.5cm]
	{\large Dualen Hochschule Baden-Württemberg Karlsruhe}\\[0.5cm]
	{\large von}\\[0.5cm]
	{\large\bfseries \Autor}\\[1cm]
	{\large Abgabedatum \AbgabeDatum}
	\vfill
\end{center}
\begin{tabular}{l@{\hspace{2cm}}l}
	Matrikelnummer	                 	& \MatrikelNummer		\\
	Kurs			         			& \Kursbezeichnung		\\
	Bearbeitungszeitraum				& \Dauer 				\\
	Gutachter der Studienakademie	 	& \BetreuerDHBW			\\
\end{tabular}
	\end{titlepage}

	% Ab hier beginnt die Seitenzählung für den Prefix in kleinen römischen Zahlen
	% Erste Seite nach der Titelseite
	\pagenumbering{roman}
	\setcounter{page}{2}
	\ifthenelse{\boolean{DEBUG}}{}{\cleardoublepage}
	
	%—————————————————————————————————————————————————————————————————————————
	%							Inhaltsverzeichnis
	%—————————————————————————————————————————————————————————————————————————
	\setcounter{tocdepth}{1}% Allow only \chapter in ToC
	\ifthenelse{\boolean{DEBUG}}{}{\cleardoublepage}
	
	\phantomsection
	\addcontentsline{toc}{chapter}{Inhaltsverzeichnis}
	\tableofcontents           % Inhaltsverzeichnis hier ausgeben
	\ifthenelse{\boolean{DEBUG}}{\newpage}{\cleardoublepage}
	
	%—————————————————————————————————————————————————————————————————————————
	%								Inhalt
	%—————————————————————————————————————————————————————————————————————————
	%Hauptteil normale arabische Zählung, beginnend bei 1 und alten Pagestyle wiederherstellen!
	\pagenumbering{arabic} 
	\setcounter{page}{1}
	\pagestyle{headings} %<- Pagestyle reset 
	
	\chapter{Domain Driven Design}
Das Konzept der Ubiquitous Language ist einer der Grundpfeiler im Domain Driven Design: In der Kommunikation mit dem Kunden und im gesamten Quellcode des Projekts sollen hierbei einheitliche Begriffe verwendet werden.

    \section{Analyse der Ubiquitous Language}
    Ein Manga besitzt einen Namen und einen Publisher, wobei ein Publisher eine Vielzahl an Mangas herausbringen kann. Zudem besitzt ein Manga einen Author und ein Genre, wobei sowohl der Author als auch das Genre von mehreren Mangas verwendet werden kann. Ein Manga hat zudem eine Liste (episodeList), in welcher mehrere Episoden sein können. Diese Episoden besitzen außerdem eine ratingList, welche Bewertungen für eine spezifische Episode speichert.
    
    \section{Analyse und Begründung der verwendeten Muster}
    Im nächsten Schritt sollen die verwendeten Muster und deren Nutzen im Projekt erläutert werden.
    
        \subsection{Value Objects}
        Das Value Object ist ein in der Softwareentwicklung eingesetztes Entwurfsmuster. Wertobjekte sind unveränderbare Objekte, die einen speziellen Wert repräsentieren. Soll der Wert geändert werden, so muss ein neues Objekt generiert werden. Sei sind zudem nur über die Werte ihrer Attribute definiert und besitzen keine eigene Identität. Zwei Wertobjekte sind identisch, wenn alle Werte ihrer Attribute identisch sind.
        
        \subsection{Entities}
        Entities sind die Grundobjekte, welche durch ihre eigene Identität definiert werden und nicht durch ihre Attribute. In diesem Projekt sind dies die folgenden Domainobjekte:
        \begin{itemize}
            \item Country 
            \item Manga
            \item Publisher
            \item Epiosde
            \item Author
            \item Rating
        \end{itemize}
        
        \subsection{Aggregates}
        In diesem Objekt sind keine Aggregates vorhanden, da keine Entitätsgruppen gebildet werden müssen. Im Normalfall werden in Aggregates Entitäten in einer logischen Gruppe zusammengefasst.

        \subsection{Repositories}
        Die Repositories dienen als Schnittstelle zur Persistierungslogik. Hierbei können die Entitäten mit Hilfe dieser Logik persistiert oder aus einer persistierten Zustand zurück in die Anwendung geladen werden.
        
        \subsection{Application Services}
        Die Application Services werden verwendet, da die Use-Cases über den Bereich einer einzelnen Entität hinausgehen. Somit wurden diese Operationen in bestimmte Appliciation Services ausgelagert.
	\ifthenelse{\boolean{DEBUG}}{}{\cleardoublepage}
	
	%%%%%%%%%%%%%%%%%%%%%%%%%%%%%%%%%%%%%%%%%%%%%%%%%%%%%%%%%%%%%%%%%%%%%%%%%%%%%%%
%% Descr:       Vorlage für Berichte der DHBW-Karlsruhe
%% Author:      Prof. Dr. Jürgen Vollmer, vollmer@dhbw-karlsruhe.de
%% Modified :	Nico Holzhäuser, TINF19B4
%% -*- coding: utf-8 -*-
%%%%%%%%%%%%%%%%%%%%%%%%%%%%%%%%%%%%%%%%%%%%%%%%%%%%%%%%%%%%%%%%%%%%%%%%%%%%%%%

\chapter{Clean Architecture}

\section{Schichtenarchitektur}

\subsection{Planung}

\subsection{Entscheidung anhand von Kriterien}
	\ifthenelse{\boolean{DEBUG}}{}{\cleardoublepage}
	
	%%%%%%%%%%%%%%%%%%%%%%%%%%%%%%%%%%%%%%%%%%%%%%%%%%%%%%%%%%%%%%%%%%%%%%%%%%%%%%%
%% Descr:       Vorlage für Berichte der DHBW-Karlsruhe
%% Author:      Prof. Dr. Jürgen Vollmer, vollmer@dhbw-karlsruhe.de
%% Modified :	Nico Holzhäuser, TINF19B4
%% -*- coding: utf-8 -*-
%%%%%%%%%%%%%%%%%%%%%%%%%%%%%%%%%%%%%%%%%%%%%%%%%%%%%%%%%%%%%%%%%%%%%%%%%%%%%%%

\chapter{Programming Principles}

\section{SOLID}
\subsection{Analyse}
\subsection{Begründung}

\section{GRASP (insb. Kopplung/Kohäsion)}
\subsection{Analyse}
\subsection{Begründung}

\section{DRY}
\subsection{Analyse}
\subsection{Begründung}
	\ifthenelse{\boolean{DEBUG}}{}{\cleardoublepage}
	
	\chapter{Refactoring}
In diesem Kapitel sollen die Refactorings betrachtet werden.
    \section{Identifizieren von Codesmells}
    Hierbei wurde das Sonar Lint Plugin verwendet, um die bestehenden Code-Smells des Projektes anzuzeigen. Dabei wurden die gefundenen Fehler mit den RISC Nummern abgeglichen, um einen Lösungsansatz zu bekommen.
        \subsection{Codesmell 1}
        \begin{figure}[h]
        	\centering
        	%\includegraphics[height=150px]{./appendix/filesONLYForAppendix/Images/lion.png}
        	\shadowimage[height=80px]{./zfiles/Bilder/CodeSmell.png}
        	\caption{Codesmell 1 - Unnötige Imports löschen}
        	\label{codesmell_1}
        \end{figure}
        Der erste Codesmell befasst sich mit den bei der Entwicklung importierten Bestandteilen. Bei der Entwicklung werden oft verschiedene Packages importiert, die wegen Codeänderungen am Ende gar nicht verwendet werden. Dieser Code-Smell mit der Kennung \hk{java:S1700} soll nun behoben werden. 
        \subsubsection{Begründung}
        Diese Imports stellen ein Sicherheitsrisiko für die Anwendung dar und tragen außerdem nicht zur besseren Lesbarkeit des Codes bei, deshalb sollen sie entfernt werden.
        \subsubsection{Fix}
        Diese Imports müssen schlichtweg aus den implementierten Klassen entfernt werden. Dies kann manuell oder mit dem Refactoring Tool der IDE umgesetzt werden. Dieses Refactoring wurde im Commit \hk{70c0170a518290b3449e383bd04e1c7e99e740fd} umgesetzt.
        
        \subsection{Codesmell 2}
        \begin{figure}[h]
        	\centering
        	%\includegraphics[height=150px]{./appendix/filesONLYForAppendix/Images/lion.png}
        	\shadowimage[height=90px]{./zfiles/Bilder/CodeSmell2.png}
        	\caption{Codesmell 2 - Variablennamen gleich dem Klassennamen}
        	\label{codesmell_2}
        \end{figure}
        Der zweite Codesmell bezieht sich auf die Namen von Variablen in Klassen. Hierbei hieß eine Variable wie die Klasse, in der sie verwendet wurde. Das ist nach java:S17000 ein schwerwiegender Codesmell
        \subsubsection{Begründung}
        Wenn das Attribut wie die Klasse heißt, verwirrt das den Entwicker und sollte neu benannt werden.
        \subsubsection{Fix}
        Umbenennung des Attributes wurde in Commit \hk{d4bb6279cfc335f8967256a795db978ae6ad6930} umgesetzt.
	\ifthenelse{\boolean{DEBUG}}{}{\cleardoublepage}
	
	\ifthenelse{\boolean{TRENNSTRICHE}}{\addtocontents{toc}{\protect\mbox{}\protect\hrulefill\par}}{}
	\chapter{Entwurfsmuster}
Im Projekt wurde das Bridge Pattern verwendet. Hierbei werden die Interfaces, welche in der Domain-Schicht entwickelt wurden mit Hilfe des Persistierungsformeworks Hibernate implementiert. Dabei wird die Implementierung von der Abstraktion entkoppelt. Dies hat zur Folge, dass die Persistierungslogik leicht ausgetauscht werden kann, wenn die neue Persistierungslogik die definierten Operationen besitzt.

    \section{UML Vorher}
    \begin{figure}[h]
    	\centering
    	%\includegraphics[height=150px]{./appendix/filesONLYForAppendix/Images/lion.png}
    	\shadowimage[height=50px]{./zfiles/Bilder/before.png}
    	\caption{Vor der Benutzung des Bridge Patterns}
    	\label{before}
    \end{figure}
    
    \section{UML Nachher}
    \begin{figure}[h]
    	\centering
    	%\includegraphics[height=150px]{./appendix/filesONLYForAppendix/Images/lion.png}
    	\shadowimage[height=150px]{./zfiles/Bilder/after.png}
    	\caption{Nach der Benutzung des Bridge Patterns}
    	\label{before}
    \end{figure}
	%Trennlinie fürs ToC beim letzten Kapitel!
	
	\ifthenelse{\boolean{DEBUG}}{}{\cleardoublepage}
	
\end{document}